% =====================================================================
% MODÉLISATION UML — ADMP AI STUDIO
% Inclus dans conception_architecture.tex (section 5.1 → 5.1.3)
% =====================================================================

\subsection{Modélisation UML de la plateforme ADMP~AI~Studio}
\label{subsec:uml-admp}

Les tableaux de synthèse de la section~\ref{subsec:avant-apres} décrivent la plateforme sur le plan fonctionnel. Cette section formalise l'analyse par une modélisation UML conforme aux bonnes pratiques du génie logiciel~: vue fonctionnelle (cas d'utilisation), vue comportementale (séquences) et vue structurelle (classes).

% =====================================================================
\subsubsection{Diagrammes de cas d'utilisation}
\label{subsubsec:uc-diagrammes}

% ---- Styles UML locaux ----
\tikzset{
  ucact/.style={rectangle, rounded corners=3pt,
                fill=accenturedark, text=white, font=\scriptsize\bfseries,
                minimum width=2.1cm, minimum height=0.55cm,
                align=center, inner sep=4pt},
  ucuse/.style={ellipse, draw=accenturepurple, line width=0.9pt,
                fill=accenturelightpurple!50, text=accenturedark,
                font=\tiny\bfseries, align=center,
                minimum width=2.8cm, minimum height=0.72cm, inner sep=2pt},
  uclin/.style={draw=accenturedark!55, line width=0.7pt},
  uczone/.style={fill=accenturedark!7, rounded corners=4pt},
  umlch/.style={fill=accenturepurple, text=white, font=\scriptsize\bfseries,
                minimum width=3.6cm, minimum height=0.6cm,
                align=center, inner sep=3pt, draw=accenturedark},
  umlcb/.style={fill=accenturegray!50, text=accenturedark, font=\tiny,
                minimum width=3.6cm, align=left, inner sep=4pt,
                draw=accenturedark, line width=0.5pt},
  umlrel/.style={draw=accenturedark!70, ->, >=open triangle 45},
  umlinher/.style={draw=accenturedark, ->, >=open triangle 60, line width=1pt},
  umlcompo/.style={draw=accenturedark, ->, >=diamond, line width=0.8pt},
  umlaggr/.style={draw=accenturedark, ->, >=open diamond, line width=0.8pt}
}

% =====================================================================
\paragraph{Version~1 --- Vue globale de la plateforme ADMP~AI~Studio}

\begin{figure}[H]
\centering
\resizebox{\textwidth}{!}{%
\begin{tikzpicture}[>=Latex]

%---- Frontière système ----
\fill[accenturegray!20, rounded corners=8pt] (0,0) rectangle (13,15.5);
\draw[accenturedark!65, line width=1.8pt, rounded corners=8pt] (0,0) rectangle (13,15.5);
\node[font=\small\bfseries, text=accenturedark] at (6.5,15.15)
  {$\langle\langle$système$\rangle\rangle$~~ADMP~AI~Studio};

%---- Zones internes ----
\fill[uczone] (0.3,13.8) rectangle (12.7,14.9);
\node[font=\tiny\bfseries, text=accenturedark!55, anchor=north west] at (0.4,14.85)
  {Administration \& Configuration};

\fill[uczone] (0.3,11.9) rectangle (12.7,13.5);
\node[font=\tiny\bfseries, text=accenturedark!55, anchor=north west] at (0.4,13.45)
  {Project Management \& Data Scope};

\fill[uczone] (0.3,6.2) rectangle (12.7,11.6);
\node[font=\tiny\bfseries, text=accenturedark!55, anchor=north west] at (0.4,11.55)
  {Extract · Profiling · Transform};

\fill[uczone] (0.3,3.5) rectangle (12.7,5.9);
\node[font=\tiny\bfseries, text=accenturedark!55, anchor=north west] at (0.4,5.85)
  {Validation \& Load};

\fill[uczone] (0.3,0.3) rectangle (12.7,3.2);
\node[font=\tiny\bfseries, text=accenturedark!55, anchor=north west] at (0.4,3.15)
  {Cutover \& Livraison};

%---- Cas d'utilisation ----
% Administration
\node[ucuse] (uc01) at (3.3,14.3)  {UC-01\\Gérer utilisateurs \& rôles};
\node[ucuse] (uc02) at (9.7,14.3)  {UC-02\\Configurer la plateforme};

% Project Management
\node[ucuse] (uc03) at (3.3,12.7)  {UC-03\\Créer un projet\\de migration};
\node[ucuse] (uc04) at (9.7,12.7)  {UC-04\\Définir le périmètre\\SDS / TDS / IDS};

% Extract · Profiling · Transform
\node[ucuse] (uc05) at (2.5,10.7)  {UC-05\\Importer \& profiler\\les données sources};
\node[ucuse] (uc06) at (6.5,10.7)  {UC-06\\Rédiger les règles\\RC et DP};
\node[ucuse] (uc07) at (10.5,10.7) {UC-07\\Configurer les\\transcodifications};
\node[ucuse] (uc08) at (2.5, 8.5)  {UC-08\\Générer le bundle\\de migration};
\node[ucuse] (uc09) at (6.5, 8.5)  {UC-09\\Réviser \& approuver\\les règles};
\node[ucuse] (uc10) at (10.5, 8.5) {UC-10\\Consulter tableaux\\de bord \& rapports};
\node[ucuse] (uc11) at (2.5, 6.7)  {UC-11\\Gérer les clés\\de décision (KD)};
\node[ucuse] (uc12) at (9.7, 6.7)  {UC-12\\Piloter les questions\\et réponses};

% Validation & Load
\node[ucuse] (uc13) at (3.3, 4.7)  {UC-13\\Exécuter les contrôles\\d'intégrité};
\node[ucuse] (uc14) at (9.7, 4.7)  {UC-14\\Réconcilier\\les données};

% Cutover
\node[ucuse] (uc15) at (3.3, 1.9)  {UC-15\\Planifier le\\basculement};
\node[ucuse] (uc16) at (9.7, 1.9)  {UC-16\\Exécuter le cutover\\et Go-Live};

%---- Acteurs (gauche) ----
\node[ucact] (adm) at (-2.8, 14.3) {Administrator};
\node[ucact] (md)  at (-2.8, 10.0) {Migration\\Designer};
\node[ucact] (da)  at (-2.8,  7.0) {Data Analyst};

%---- Acteurs (droite) ----
\node[ucact] (te)  at (15.8, 11.5) {Test Engineer};
\node[ucact] (br)  at (15.8,  8.5) {Business\\Reviewer};
\node[ucact] (cm)  at (15.8,  2.5) {Cutover\\Manager};
\node[ucact] (cl)  at (15.8,  0.0) {Client};

%---- Associations ----
\draw[uclin] (adm) -- (uc01);
\draw[uclin] (adm) -- (uc02);
\draw[uclin] (md)  -- (uc03);
\draw[uclin] (md)  -- (uc04);
\draw[uclin] (md)  -- (uc08);
\draw[uclin] (md)  -- (uc10);
\draw[uclin] (da)  -- (uc05);
\draw[uclin] (da)  -- (uc06);
\draw[uclin] (da)  -- (uc07);
\draw[uclin] (da)  -- (uc11);
\draw[uclin] (da)  -- (uc12);
\draw[uclin] (te)  -- (uc13);
\draw[uclin] (te)  -- (uc14);
\draw[uclin] (te)  -- (uc10);
\draw[uclin] (br)  -- (uc09);
\draw[uclin] (br)  -- (uc12);
\draw[uclin] (br)  -- (uc14);
\draw[uclin] (cm)  -- (uc15);
\draw[uclin] (cm)  -- (uc16);
\draw[uclin] (cl)  -- (uc10);
\draw[uclin] (cl)  -- (uc16);

\end{tikzpicture}%
}
\caption{Diagramme de cas d'utilisation --- Vue globale de la plateforme ADMP~AI~Studio}
\label{fig:uc-global}
\end{figure}

% =====================================================================
\paragraph{Version~2 --- Module de génération Databricks (contribution PFE)}

\begin{figure}[H]
\centering
\begin{tikzpicture}[>=Latex]

%---- Frontière module ----
\fill[accenturelightpurple!30, rounded corners=8pt] (0,0) rectangle (10,9.5);
\draw[accenturepurple, line width=1.8pt, rounded corners=8pt] (0,0) rectangle (10,9.5);
\node[font=\small\bfseries, text=accenturepurple] at (5, 9.15)
  {$\langle\langle$module$\rangle\rangle$~~Package Generation --- Databricks};

%---- Zones internes ----
\fill[accenturepurple!10, rounded corners=3pt] (0.3,7.0) rectangle (9.7,8.7);
\node[font=\tiny\bfseries, text=accenturepurple!70, anchor=north west] at (0.4,8.65)
  {Configuration \& Compilation};

\fill[accenturepurple!10, rounded corners=3pt] (0.3,3.8) rectangle (9.7,6.7);
\node[font=\tiny\bfseries, text=accenturepurple!70, anchor=north west] at (0.4,6.65)
  {Déploiement \& Exécution};

\fill[accenturepurple!10, rounded corners=3pt] (0.3,0.3) rectangle (9.7,3.5);
\node[font=\tiny\bfseries, text=accenturepurple!70, anchor=north west] at (0.4,3.45)
  {Validation \& Reporting};

%---- Cas d'utilisation ----
\node[ucuse] (ucd01) at (2.8,7.85) {UC-DB01\\Configurer le format\\Databricks};
\node[ucuse] (ucd02) at (7.2,7.85) {UC-DB02\\Lancer la compilation\\du bundle};
\node[ucuse] (ucd03) at (2.8,5.35) {UC-DB03\\Visualiser le code\\PySpark généré};
\node[ucuse] (ucd04) at (7.2,5.35) {UC-DB04\\Déployer et exécuter\\les unités PySpark};
\node[ucuse] (ucd05) at (2.8,2.0)  {UC-DB05\\Consulter les rapports\\d'exécution Excel};
\node[ucuse] (ucd06) at (7.2,2.0)  {UC-DB06\\Diagnostiquer et corriger\\les erreurs SQL};

%---- Acteurs ----
\node[ucact] (dbt-md)  at (-2.6,7.85) {Migration\\Designer};
\node[ucact] (dbt-da)  at (-2.6,4.35) {Data Analyst};
\node[ucact] (dbt-te)  at (12.6, 3.0) {Test Engineer};

%---- Associations ----
\draw[uclin] (dbt-md) -- (ucd01);
\draw[uclin] (dbt-md) -- (ucd02);
\draw[uclin] (dbt-da) -- (ucd03);
\draw[uclin] (dbt-da) -- (ucd04);
\draw[uclin] (dbt-da) -- (ucd06);
\draw[uclin] (dbt-te) -- (ucd05);
\draw[uclin] (dbt-te) -- (ucd06);

%---- Relation include ----
\draw[draw=accenturepurple!70, dashed, ->, >=Stealth]
  (ucd02) -- (ucd03) node[midway, right, font=\tiny\itshape, text=accenturepurple!70]
  {$\langle\langle$include$\rangle\rangle$};

\end{tikzpicture}
\caption{Diagramme de cas d'utilisation --- Module de génération Databricks (contribution PFE)}
\label{fig:uc-databricks}
\end{figure}

% =====================================================================
\subsubsection{Tableau récapitulatif des cas d'utilisation}
\label{subsubsec:uc-table}

\begin{table}[H]
\centering
\caption{Récapitulatif des cas d'utilisation de la plateforme ADMP~AI~Studio}
\label{tab:uc-recap}
\begin{tabularx}{\textwidth}{llX}
\toprule
\textbf{Acteur(s)} & \textbf{Cas d'utilisation} & \textbf{Description} \\
\midrule
Administrator &
  UC-01 Gérer utilisateurs \& rôles &
  Créer, modifier, désactiver des comptes utilisateurs~; assigner des rôles et des permissions par projet \\
\addlinespace
Administrator &
  UC-02 Configurer la plateforme &
  Paramétrer les providers LLM (\texttt{llm\_credentials}), les préférences système et les invites contextuelles \\
\addlinespace
Migration Designer &
  UC-03 Créer un projet de migration &
  Initialiser un nouveau projet~: Industry, Target, Source, phase courante, membres (\texttt{projects}, \texttt{project\_members}) \\
\addlinespace
Migration Designer, Data Analyst &
  UC-04 Définir le périmètre &
  Rédiger et publier les documents SDS, TDS et IDS (\texttt{datascope\_documents}, \texttt{datascope\_tables}) \\
\addlinespace
Data Analyst &
  UC-05 Importer \& profiler les sources &
  Importer les fichiers plats~; exécuter le profilage des champs (\texttt{data\_extractions}, \texttt{source\_audit\_*}) \\
\addlinespace
Data Analyst &
  UC-06 Rédiger les règles RC / DP &
  Rédiger les instructions Record Creation et Data Population avec checkout/check-in (\texttt{migration\_units}) \\
\addlinespace
Data Analyst &
  UC-07 Configurer les transcodifications &
  Créer et versionner les tables de transco~; définir colonnes et mappings (\texttt{transform\_transco\_*}) \\
\addlinespace
Migration Designer &
  UC-08 Générer le bundle &
  Déclencher la compilation et récupérer le package \texttt{.whl} pour Databricks (\texttt{configuration\_requests}) \\
\addlinespace
Business Reviewer &
  UC-09 Réviser \& approuver &
  Valider les règles de transformation et data mapping~; approuver les versions publiées \\
\addlinespace
Migration Designer, Test Engineer, Client &
  UC-10 Consulter tableaux \& rapports &
  Accéder aux tableaux de bord de progression, rapports de management hebdomadaires (\texttt{management\_reports}) \\
\addlinespace
Data Analyst &
  UC-11 Gérer les clés de décision &
  Enregistrer, commenter et clore les décisions de conception (\texttt{key\_decisions}) \\
\addlinespace
Data Analyst, Business Reviewer &
  UC-12 Piloter les questions/réponses &
  Poser et répondre aux questions inter-équipes (\texttt{questions}, \texttt{question\_comments}) \\
\addlinespace
Test Engineer &
  UC-13 Exécuter les contrôles d'intégrité &
  Lancer les vérifications référentielles et consulter les résultats (\texttt{integrity\_checks}, \texttt{validation\_reports}) \\
\addlinespace
Test Engineer, Business Reviewer &
  UC-14 Réconcilier les données &
  Comparer source et cible aux niveaux L1 (comptage), L2 (champ), L3 (règles métier) (\texttt{reconciliation\_*}) \\
\addlinespace
Cutover Manager &
  UC-15 Planifier le basculement &
  Créer le plan de cutover, définir tâches, jalons et approbateurs (\texttt{cutover\_plans}, \texttt{cutover\_tasks}) \\
\addlinespace
Cutover Manager, Client &
  UC-16 Exécuter le cutover \& Go-Live &
  Démarrer l'exécution, suivre les tâches, gérer les incidents (\texttt{cutover\_executions}, \texttt{cutover\_issues}) \\
\addlinespace
Migration Designer &
  UC-DB01 Configurer format Databricks &
  Sélectionner Databricks parmi les 8 formats disponibles et paramétrer catalog / schema \\
\addlinespace
Migration Designer &
  UC-DB02 Lancer la compilation &
  Soumettre la requête de compilation~; le compilateur génère le bundle \texttt{.whl} \\
\addlinespace
Data Analyst &
  UC-DB03 Visualiser le code PySpark &
  Inspecter les scripts \texttt{mig\_*.py}, les requêtes SQL et le framework généré \\
\addlinespace
Data Analyst &
  UC-DB04 Déployer et exécuter &
  Déployer la librairie sur le cluster Databricks~; déclencher les jobs PySpark \\
\addlinespace
Test Engineer &
  UC-DB05 Consulter les rapports Excel &
  Récupérer et analyser les rapports d'intégrité exportés depuis le volume Unity Catalog \\
\addlinespace
Data Analyst, Test Engineer &
  UC-DB06 Diagnostiquer les erreurs &
  Analyser les erreurs SQL Spark~; corriger les règles RC/DP puis regénérer \\
\bottomrule
\end{tabularx}
\end{table}

% =====================================================================
\subsubsection{Description détaillée des cas d'utilisation principaux}
\label{subsubsec:uc-details}

\newcommand{\ficheuc}[7]{%
\begin{tcolorbox}[
  colback=accenturegray!40, colframe=accenturepurple,
  title={\textbf{#1}}, fonttitle=\small\bfseries,
  boxrule=0.8pt, arc=4pt, left=6pt, right=6pt, top=4pt, bottom=4pt
]
\begin{tabularx}{\linewidth}{lX}
\textbf{Acteur(s)} & #2 \\
\textbf{Objectif}  & #3 \\
\textbf{Pré-conditions} & #4 \\
\textbf{Post-conditions} & #5 \\
\end{tabularx}
\medskip\par
\textbf{Scénario nominal~:}\vspace{2pt}
\begin{enumerate}[leftmargin=*, label=\textbf{\arabic*.}, nosep]
#6
\end{enumerate}
\ifx&#7&\else
\medskip\par\textbf{Scénarios alternatifs / extensions~:}\par#7
\fi
\end{tcolorbox}
\vspace{6pt}
}

% ---- UC-03 ----
\ficheuc{UC-03 --- Créer un projet de migration}%
{Migration Designer}%
{Ouvrir un nouveau projet ADMP sur la plateforme ADMP~AI~Studio en définissant ses métadonnées, son équipe et son périmètre initial.}%
{L'utilisateur est authentifié et possède le rôle \textit{Migration Designer}. Au moins un système source et un système cible sont configurés dans la plateforme.}%
{Un enregistrement est créé dans \texttt{projects} (statut \texttt{active}). L'utilisateur est ajouté comme membre dans \texttt{project\_members} avec le rôle \textit{owner}. La phase \texttt{project\_management} est initialisée dans \texttt{project\_phases}.}%
{%
\item L'utilisateur clique sur « Nouveau projet » depuis le tableau de bord.
\item Il renseigne les métadonnées~: nom, client, Industry, système source, système cible, dates prévisionnelles.
\item Il définit le Data Scope initial en indiquant les périmètres SDS, TDS, IDS concernés.
\item Il ajoute les membres de l'équipe et leur attribue un rôle (\textit{Migration Designer}, \textit{Data Analyst}, etc.).
\item Il soumet le formulaire~; la plateforme crée le projet et redirige vers le tableau de bord projet.
\item Les huit phases de migration apparaissent avec le statut \texttt{not\_started}.
}%
{2a. Les champs obligatoires sont manquants~: la plateforme affiche les erreurs de validation et empêche la soumission. / 3a. Le système source n'est pas encore référencé~: l'utilisateur peut le créer à la volée depuis le formulaire.}

% ---- UC-06 ----
\ficheuc{UC-06 --- Rédiger les règles RC et DP}%
{Data Analyst}%
{Rédiger ou modifier les instructions Record Creation et Data Population d'une unité de migration, avec verrouillage de section pour éviter les conflits d'édition concurrents.}%
{Le projet est en phase \textit{Transform}. L'unité de migration (\texttt{migration\_units}) existe et sa section RC (ou DP) n'est pas actuellement verrouillée par un autre utilisateur (\texttt{*\_locked\_by IS NULL}).}%
{Le contenu RC/DP est enregistré dans \texttt{migration\_units}. Une nouvelle entrée est créée dans \texttt{migration\_unit\_versions} (section = \texttt{record\_creation} ou \texttt{data\_population}, statut \texttt{draft}). La section est déverrouillée (\texttt{*\_locked\_by = NULL}).}%
{%
\item L'utilisateur sélectionne l'unité de migration dans la liste.
\item Il clique sur « Checkout » pour la section RC (ou DP).
\item La plateforme enregistre \texttt{*\_locked\_by = userId}, \texttt{*\_locked\_at = NOW()}, et un snapshot de l'état courant dans \texttt{*\_checkout\_baseline}.
\item L'utilisateur rédige les instructions dans l'éditeur (IF/ELSE, callWriteFile, callDataPopulate, assignations).
\item Il clique sur « Check-in »~; la plateforme valide la syntaxe et crée une version \texttt{draft}.
\item Le verrou est libéré (\texttt{*\_locked\_by = NULL}).
}%
{4a. Un autre utilisateur tente un checkout pendant que la section est verrouillée~: la plateforme retourne une erreur 409 avec le nom de l'utilisateur propriétaire du verrou. / 5a. L'utilisateur clique sur « Discard » au lieu de « Check-in »~: le contenu est rétabli depuis \texttt{*\_checkout\_baseline}.}

% ---- UC-08 ----
\ficheuc{UC-08 --- Générer le bundle de migration}%
{Migration Designer}%
{Déclencher la compilation de l'ensemble des spécifications ADMP d'un projet en un bundle PySpark déployable sur Databricks.}%
{Le projet est en phase \textit{Transform}. Au moins une unité de migration possède des versions publiées pour les sections DS, RC et DP. Le format de sortie \textit{Databricks} est sélectionné.}%
{Un fichier \texttt{.whl} est généré et mis à disposition. Un enregistrement est créé dans \texttt{configuration\_requests} (statut \texttt{completed}). Les scripts \texttt{mig\_*.py}, \texttt{runner.py}, \texttt{validate.py}, \texttt{unit\_base.py} et \texttt{create\_tables.py} sont inclus dans le bundle.}%
{%
\item L'utilisateur navigue vers la section \textit{Package Generation} de la phase \textit{Transform}.
\item Il sélectionne le format de sortie « Databricks » et renseigne les paramètres~: catalog cible, schéma, cluster.
\item Il clique sur « Générer »~; la requête est envoyée au service de compilation (\texttt{configuration\_requests}).
\item Le compilateur TypeScript parse les documents RC, DP, TT, DS et construit les objets \texttt{SqlPath}.
\item La fonction \texttt{generateSelectBlock()} génère les blocs SQL \texttt{UNION ALL} pour chaque branche.
\item Le Bundle Assembler produit les fichiers Python et les packague en \texttt{.whl}.
\item La plateforme notifie l'utilisateur et propose le téléchargement du bundle.
}%
{4a. Un document RC contient une erreur de syntaxe~: le compilateur lève une exception~; la plateforme affiche l'erreur avec la ligne incriminée. / 6a. Aucune unité publiée n'est trouvée~: la plateforme affiche un avertissement et empêche la compilation.}

% ---- UC-13 ----
\ficheuc{UC-13 --- Exécuter les contrôles d'intégrité}%
{Test Engineer}%
{Lancer les vérifications d'intégrité référentielle sur les données migrées et produire un rapport Excel consolidé par unité.}%
{Le bundle Databricks a été déployé et au moins un job d'exécution s'est terminé avec succès. Les règles d'intégrité sont définies dans \texttt{integrity\_checks}.}%
{Les résultats sont enregistrés dans \texttt{validation\_reports}. Un fichier Excel est écrit dans le volume Unity Catalog \texttt{dbfs:/Volumes/\textit{catalog}/\textit{schema}/reports/\textit{runId}/}.}%
{%
\item L'utilisateur navigue vers la section \textit{Validation} du projet.
\item Il sélectionne le run Databricks à analyser et clique sur « Lancer les contrôles ».
\item Le job \texttt{validate.py} interroge les tables Delta cibles pour chaque contrainte définie.
\item Les résultats (comptages, violations, taux de conformité) sont calculés par unité.
\item Un rapport Excel \texttt{integrity\_report\_\textit{unitCode}.xlsx} est écrit dans le volume Unity Catalog.
\item Un rapport consolidé \texttt{global\_report.xlsx} agrège toutes les unités.
\item La plateforme affiche le tableau de résultats et met à disposition les fichiers Excel via le script \texttt{fetch\_reports.ps1}.
}%
{4a. Une violation d'intégrité critique est détectée~: le rapport la signale en rouge~; le Test Engineer ouvre un ticket dans \texttt{test\_tickets}. / 3a. Le volume Unity Catalog est inaccessible~: le job lève une exception~; le run est marqué \texttt{failed}.}

% ---- UC-15 ----
\ficheuc{UC-15 --- Planifier le basculement (Cutover)}%
{Cutover Manager}%
{Créer et structurer le plan de basculement définitif en production, incluant les tâches séquencées, les jalons, les approbateurs et le plan de rollback.}%
{Le projet est en phase \textit{Cutover}. Les réconciliations L1, L2 et L3 ont été validées. Au moins un Cutover Manager est désigné dans \texttt{project\_members}.}%
{Un enregistrement est créé dans \texttt{cutover\_plans} (statut \texttt{draft}). Les tâches sont créées dans \texttt{cutover\_tasks} avec leurs dépendances. Les jalons sont définis dans \texttt{cutover\_milestones}.}%
{%
\item Le Cutover Manager crée un nouveau plan depuis la section \textit{Cutover} du projet.
\item Il renseigne les métadonnées~: date cible, durée estimée, environnement, plan de rollback, plan de communication.
\item Il décompose le plan en tâches (import, arrêt des interfaces, chargement final, validation, bascule DNS, etc.) avec leurs durées et dépendances.
\item Il marque les tâches sur le chemin critique (\texttt{critical\_path = true}).
\item Il définit les jalons (\texttt{cutover\_milestones}) et les approbateurs (\texttt{cutover\_approvers}).
\item Il soumet le plan pour approbation~; les approbateurs reçoivent une notification.
\item Après approbation, le plan passe au statut \texttt{approved}.
}%
{6a. Un approbateur rejette le plan~: il ajoute un commentaire et le plan reste en statut \texttt{review}. / 3a. Un modèle de plan existe (\texttt{is\_template = true})~: le Cutover Manager peut partir du modèle et l'adapter.}

% =====================================================================
\subsubsection{Diagrammes de séquence}
\label{subsubsec:sequence-diagrams}

% =====================================================================
\paragraph{SD-01 --- Génération du bundle Databricks}

\begin{figure}[H]
\centering
\begin{tikzpicture}[>=Latex, thick,
  sactor/.style={rectangle, rounded corners=3pt, fill=accenturedark, text=white,
                 font=\bfseries\scriptsize, minimum width=2.4cm,
                 minimum height=0.65cm, align=center},
  smsg/.style={font=\scriptsize, text=accenturedark, align=center},
  sret/.style={font=\scriptsize, text=accenturedark!65, align=center}
]
  %---- Acteurs ----
  \node[sactor] (ui)   at (0,   0) {ADMP\\Studio UI};
  \node[sactor] (api)  at (4,   0) {Compilation\\Service (API)};
  \node[sactor] (comp) at (8,   0) {Databricks\\Compiler (TS)};
  \node[sactor] (dbk)  at (12,  0) {Databricks\\Runtime};

  %---- Lignes de vie ----
  \foreach \x in {0, 4, 8, 12} {
    \draw[accenturedark!25, dashed] (\x,-0.33) -- (\x,-11.5);
  }

  %---- Messages ----
  \draw[->] (0,-0.9)  -- (4,-0.9)
    node[midway, above, smsg] {1. POST /codegen/bundle (projectId, format=databricks)};
  \draw[->] (4,-1.7)  -- (8,-1.7)
    node[midway, above, smsg] {2. compile(units, transcos, params)};
  \draw[->, accenturepurple] (7.7,-2.5) arc[start angle=0, end angle=-180, radius=0.25];
  \node[smsg, text=accenturepurple] at (8,-2.9) {3. parseDocuments(RC, DP, TT, DS)};
  \draw[->, accenturepurple] (7.7,-3.8) arc[start angle=0, end angle=-180, radius=0.25];
  \node[smsg, text=accenturepurple] at (8,-4.2) {4. extractPaths(ast) × N unités};
  \draw[->, accenturepurple] (7.7,-5.1) arc[start angle=0, end angle=-180, radius=0.25];
  \node[smsg, text=accenturepurple] at (8,-5.5) {5. generateSelectBlock(path) × N branches};
  \draw[->, accenturepurple] (7.7,-6.4) arc[start angle=0, end angle=-180, radius=0.25];
  \node[smsg, text=accenturepurple] at (8,-6.8) {6. emitFramework() + emitIntegrityChecks()};
  \draw[->] (8,-7.7) -- (4,-7.7)
    node[midway, above, sret] {7. bundle (.whl)};
  \draw[->] (4,-8.5) -- (12,-8.5)
    node[midway, above, smsg] {8. deployLibrary(bundle) via Databricks REST API};
  \draw[->] (4,-9.3) -- (0,-9.3)
    node[midway, above, sret] {9. bundleUrl, runId};
  \draw[->] (0,-10.2) -- (12,-10.2)
    node[midway, above, smsg] {10. triggerJob(runner.py, catalog, schema)};
  \draw[<-] (0,-11.0) -- (12,-11.0)
    node[midway, above, sret] {11. jobResult (status, records, reportUrl)};
\end{tikzpicture}
\caption{Diagramme de séquence SD-01 --- Génération et exécution d'un bundle Databricks}
\label{fig:sd-bundle}
\end{figure}

\noindent\textbf{Description~:} Le Migration Designer déclenche la compilation depuis l'interface ADMP~AI~Studio (étape~1). L'API de compilation invoque le compilateur TypeScript (étape~2) qui parse les documents de spécification, extrait les \texttt{SqlPath} et génère les blocs SQL \texttt{UNION ALL} (étapes~3--6). Le bundle \texttt{.whl} est retourné à l'API (étape~7), déployé sur le cluster Databricks via l'API REST (étape~8), puis le job \texttt{runner.py} est déclenché (étape~10). Le résultat est communiqué à l'interface (étape~11).

% =====================================================================
\paragraph{SD-02 --- Checkout / Check-in d'une section RC}

\begin{figure}[H]
\centering
\begin{tikzpicture}[>=Latex, thick,
  sactor/.style={rectangle, rounded corners=3pt, fill=accenturedark, text=white,
                 font=\bfseries\scriptsize, minimum width=2.4cm,
                 minimum height=0.65cm, align=center},
  smsg/.style={font=\scriptsize, text=accenturedark, align=center},
  sret/.style={font=\scriptsize, text=accenturedark!65, align=center},
  sdb/.style={cylinder, shape border rotate=90, draw=accenturedark,
              fill=accenturelightpurple!40, font=\scriptsize\bfseries,
              minimum width=2cm, minimum height=0.65cm, align=center}
]
  \node[sactor]  (da)   at (0, 0)  {Data Analyst};
  \node[sactor]  (ui)   at (4, 0)  {ADMP\\Studio UI};
  \node[sactor]  (api)  at (8, 0)  {Unit Version\\Service (API)};
  \node[sdb]     (db)   at (12, 0) {PostgreSQL\\(migration\_units)};

  \foreach \x in {0, 4, 8, 12} {
    \draw[accenturedark!25, dashed] (\x,-0.38) -- (\x,-10.5);
  }

  \draw[->] (0,-0.9) -- (4,-0.9) node[midway, above, smsg] {1. Clic « Checkout RC »};
  \draw[->] (4,-1.7) -- (8,-1.7) node[midway, above, smsg] {2. POST /units/:id/checkout (section=rc)};
  \draw[->] (8,-2.5) -- (12,-2.5) node[midway, above, smsg] {3. SELECT locked\_by, checkout\_baseline};
  \draw[<-] (8,-3.3) -- (12,-3.3) node[midway, above, sret] {4. locked\_by = NULL → disponible};
  \draw[->] (8,-4.1) -- (12,-4.1) node[midway, above, smsg] {5. UPDATE record\_creation\_locked\_by=userId,\\locked\_at=NOW(), checkout\_baseline=content};
  \draw[<-] (4,-5.0) -- (8,-5.0) node[midway, above, sret] {6. 200 OK, éditeur déverrouillé};
  \draw[->] (0,-5.8) -- (4,-5.8) node[midway, above, smsg] {7. Édition du RC (IF/ELSE, callWriteFile)};
  \draw[->] (0,-6.6) -- (4,-6.6) node[midway, above, smsg] {8. Clic « Check-in »};
  \draw[->] (4,-7.4) -- (8,-7.4) node[midway, above, smsg] {9. POST /units/:id/checkin (section=rc, content)};
  \draw[->] (8,-8.2) -- (12,-8.2) node[midway, above, smsg] {10. INSERT migration\_unit\_versions\\(section=rc, version=N+1, status=draft)};
  \draw[->] (8,-9.0) -- (12,-9.0) node[midway, above, smsg] {11. UPDATE locked\_by=NULL};
  \draw[<-] (0,-9.8) -- (8,-9.8) node[midway, above, sret] {12. 200 OK, version N+1 créée};
\end{tikzpicture}
\caption{Diagramme de séquence SD-02 --- Checkout / Check-in d'une section RC}
\label{fig:sd-checkout}
\end{figure}

\noindent\textbf{Description~:} Ce diagramme illustre le mécanisme de verrouillage optimiste par section. L'utilisateur demande un checkout (étape~1--2)~; l'API vérifie que le verrou est libre (étapes~3--4), pose le verrou et sauvegarde le snapshot de l'état courant (\texttt{checkout\_baseline}, étape~5). Après édition, le check-in crée une nouvelle version \textit{draft} (étape~10) et libère le verrou (étape~11). En cas de discard, le contenu est rétabli depuis le snapshot.

% =====================================================================
\paragraph{SD-03 --- Réconciliation des données (niveaux L1 / L2 / L3)}

\begin{figure}[H]
\centering
\begin{tikzpicture}[>=Latex, thick,
  sactor/.style={rectangle, rounded corners=3pt, fill=accenturedark, text=white,
                 font=\bfseries\scriptsize, minimum width=2.4cm,
                 minimum height=0.65cm, align=center},
  smsg/.style={font=\scriptsize, text=accenturedark, align=center},
  sret/.style={font=\scriptsize, text=accenturedark!65, align=center}
]
  \node[sactor] (te)   at (0, 0)   {Test\\Engineer};
  \node[sactor] (ui)   at (4, 0)   {ADMP\\Studio UI};
  \node[sactor] (recon) at (8, 0)  {Reconciliation\\Service};
  \node[sactor] (data) at (12, 0)  {Databricks\\Delta Lake};

  \foreach \x in {0, 4, 8, 12} {
    \draw[accenturedark!25, dashed] (\x,-0.33) -- (\x,-11.0);
  }

  \draw[->] (0,-0.9)  -- (4,-0.9) node[midway, above, smsg] {1. Lancer réconciliation (config\_id)};
  \draw[->] (4,-1.7)  -- (8,-1.7) node[midway, above, smsg] {2. POST /reconciliation/execute};
  %-- L1
  \draw[->, accenturepurple] (8,-2.5) -- (12,-2.5)
    node[midway, above, smsg, text=accenturepurple] {3. COUNT(*) source vs cible (L1)};
  \draw[<-, accenturepurple] (8,-3.3) -- (12,-3.3)
    node[midway, above, sret, text=accenturepurple] {4. \{src: N, tgt: M, diff: N-M\}};
  %-- L2
  \draw[->, accenturepurple] (8,-4.1) -- (12,-4.1)
    node[midway, above, smsg, text=accenturepurple] {5. Comparer champs clés (L2)};
  \draw[<-, accenturepurple] (8,-4.9) -- (12,-4.9)
    node[midway, above, sret, text=accenturepurple] {6. \{matched, mismatched, field\_comparisons\}};
  %-- L3
  \draw[->, accenturepurple] (8,-5.7) -- (12,-5.7)
    node[midway, above, smsg, text=accenturepurple] {7. Vérifier KPI \& règles métier (L3)};
  \draw[<-, accenturepurple] (8,-6.5) -- (12,-6.5)
    node[midway, above, sret, text=accenturepurple] {8. \{kpi\_checks, overall\_score\}};
  %-- Persist
  \draw[->] (8,-7.3)  -- (8,-7.3);
  \draw[->, accenturepurple] (7.7,-7.3) arc[start angle=0, end angle=-180, radius=0.25];
  \node[smsg, text=accenturepurple] at (8,-7.7)
    {9. INSERT level1/2/3\_reconciliations};
  %-- Return
  \draw[->] (8,-8.5) -- (4,-8.5)
    node[midway, above, sret] {10. Résultats L1+L2+L3, overall\_score};
  \draw[<-] (0,-9.3) -- (4,-9.3)
    node[midway, above, sret] {11. Tableau de bord réconciliation};
  \draw[->] (0,-10.2) -- (4,-10.2)
    node[midway, above, smsg] {12. Ouvrir ticket si violation critique};
\end{tikzpicture}
\caption{Diagramme de séquence SD-03 --- Réconciliation des données (L1 / L2 / L3)}
\label{fig:sd-recon}
\end{figure}

\noindent\textbf{Description~:} Le Test Engineer lance une réconciliation tripartite. Le niveau L1 vérifie les comptages (nombre de lignes source vs cible). Le niveau L2 compare les valeurs champ à champ sur les clés primaires. Le niveau L3 vérifie les KPI métier et les règles d'agrégation. Les résultats sont persistés dans les tables \texttt{level1/2/3\_reconciliations} et un score global est calculé. En cas de violation critique, le Test Engineer ouvre un ticket de défaut (\texttt{test\_tickets}).

% =====================================================================
\subsubsection{Diagramme de classes}
\label{subsubsec:class-diagram}

\begin{figure}[H]
\centering
\resizebox{\textwidth}{!}{%
\begin{tikzpicture}[>=Latex, node distance=0.5cm and 1.2cm]

%---- Macros de classe ----
\newcommand{\umlclass}[5]{% name, x, y, attributes, methods
  \node[umlch, minimum width=3.6cm] (#1h) at (#2,#3) {#1};
  \node[umlcb, minimum width=3.6cm, below=0pt of #1h, text width=3.4cm] (#1a) {#4};
  \node[umlcb, minimum width=3.6cm, below=0pt of #1a, text width=3.4cm,
        fill=white!80] (#1m) {#5};
}

%---- Classes ----
% Project
\umlclass{Project}{0}{0}%
{+ id : UUID\\+ name : String\\+ status : String\\+ current\_phase : String\\+ source\_system : String\\+ target\_system : String}%
{+ getMembers()\\+ getPhases()\\+ getUnits()}

% User
\umlclass{User}{6.5}{0}%
{+ id : UUID\\+ email : String\\+ role : String\\+ is\_active : Boolean}%
{+ getProjects()\\+ getRoles()}

% ProjectMember
\umlclass{ProjectMember}{3.2}{-5.2}%
{+ id : UUID\\+ project\_id : UUID\\+ user\_id : UUID\\+ role : String\\+ permissions : JSONB}%
{+ checkout()\\+ hasPermission()}

% MigrationUnit
\umlclass{MigrationUnit}{-5.5}{-5.2}%
{+ id : UUID\\+ code : String\\+ rc\_locked\_by : String\\+ dp\_locked\_by : String\\+ rc\_checkout\_baseline : JSONB}%
{+ checkout(section)\\+ checkin(section)\\+ discard(section)}

% MigrationUnitVersion
\umlclass{MigrationUnitVersion}{-5.5}{-11.0}%
{+ id : UUID\\+ unit\_id : UUID\\+ version : String\\+ section : String\\+ status : String}%
{+ publish()\\+ rollback()}

% TranscoTable
\umlclass{TranscoTable}{1.5}{-11.0}%
{+ id : UUID\\+ name : String\\+ description : String\\+ is\_archived : Boolean}%
{+ getColumns()\\+ getMappings()\\+ publish()}

% Bundle
\umlclass{Bundle}{8.0}{-5.2}%
{+ run\_id : String\\+ format : String\\+ catalog : String\\+ schema : String\\+ status : String}%
{+ compile()\\+ deploy()\\+ fetchReports()}

% IntegrityCheck
\umlclass{IntegrityCheck}{8.0}{-11.0}%
{+ id : UUID\\+ constraint\_type : String\\+ table\_name : String\\+ key\_columns : String[]\\+ status : String}%
{+ execute()\\+ generateReport()}

% CutoverPlan
\umlclass{CutoverPlan}{-5.5}{-17.5}%
{+ id : UUID\\+ planned\_date : Date\\+ status : String\\+ rollback\_plan : Text\\+ critical\_path : Boolean}%
{+ addTask()\\+ approve()\\+ execute()}

% ExecutionRun
\umlclass{ExecutionRun}{1.5}{-17.5}%
{+ id : UUID\\+ run\_type : String\\+ status : String\\+ total\_records : Int\\+ failed\_records : Int}%
{+ start()\\+ getStats()\\+ getLogs()}

%---- Relations ----
% Project --- ProjectMember
\draw[umlaggr] (Projecta.south) |- ++(0,-0.5) -| (ProjectMemberh.north);
\node[font=\tiny, text=accenturedark] at (-0.4,-2.5) {1};
\node[font=\tiny, text=accenturedark] at (2.0,-2.5) {*};

% User --- ProjectMember
\draw[umlaggr] (Usera.south) |- ++(0,-0.5) -| (ProjectMemberh.north);
\node[font=\tiny, text=accenturedark] at (6.4,-2.5) {1};
\node[font=\tiny, text=accenturedark] at (4.5,-2.5) {*};

% Project --- MigrationUnit
\draw[umlcompo] (Projectm.south) -- ++(0,-0.8) -| (MigrationUnith.north);
\node[font=\tiny, text=accenturedark] at (-3.5,-2.8) {1..*};

% MigrationUnit --- MigrationUnitVersion
\draw[umlcompo] (MigrationUnita.south) -- (MigrationUnitVersionh.north);
\node[font=\tiny, text=accenturedark] at (-6.8,-8.0) {1};
\node[font=\tiny, text=accenturedark] at (-6.8,-9.0) {*};

% Project --- Bundle
\draw[umlaggr] (Projectm.east) -| (Bundleh.north);
\node[font=\tiny, text=accenturedark] at (5.5,-2.5) {*};

% Bundle --- ExecutionRun
\draw[umlcompo] (Bundlem.south) -- ++(0,-0.5) -| (ExecutionRunh.north);
\node[font=\tiny, text=accenturedark] at (6.5,-14.5) {1..*};

% MigrationUnit --- TranscoTable
\draw[umlrel, dashed] (MigrationUnita.east) -| (TranscoTableh.north);
\node[font=\tiny\itshape, text=accenturepurple] at (-1.5,-9.5) {utilise};

% Bundle --- IntegrityCheck
\draw[umlcompo] (Bundlem.south) -- ++(0,-0.3) -| (IntegrityCheckh.north);
\node[font=\tiny, text=accenturedark] at (9.5,-8.0) {1..*};

% Project --- CutoverPlan
\draw[umlcompo] (MigrationUnita.south) -- ++(0,-0.3) -| (CutoverPlanh.north);
\node[font=\tiny, text=accenturedark] at (-6.8,-14.5) {0..*};

\end{tikzpicture}%
}
\caption{Diagramme de classes --- Entités principales de la plateforme ADMP~AI~Studio}
\label{fig:class-diagram}
\end{figure}

\noindent Le diagramme de classes met en évidence les entités structurantes de la plateforme. La classe \textbf{Project} est le pivot central~: elle agrège les membres (\texttt{ProjectMember}), les unités de migration (\texttt{MigrationUnit}) et les bundles générés. La classe \textbf{MigrationUnit} porte le mécanisme de checkout/check-in~: trois paires d'attributs \texttt{*\_locked\_by} / \texttt{*\_checkout\_baseline} garantissent l'édition concurrente sans conflit sur les sections DS, RC et DP. La classe \textbf{Bundle} encapsule le cycle de compilation-déploiement-exécution et génère des \textbf{IntegrityCheck} utilisés pour la validation post-migration. Le \textbf{CutoverPlan} formalise le basculement final, constitué de tâches (\texttt{CutoverTask}), jalons (\texttt{CutoverMilestone}) et incidents (\texttt{CutoverIssue}) gérés à l'exécution.
